% answers_short.tex —— 2004 年·全国卷Ⅱ·理科 · 简短答案速查
% 由 tools/make_short_answers.py 自 answers.tex 的 \gksolhead/\gkanswer 提取，
% 并与 2004/2004_数学真题_全国2理_答案要点.md 逐题比对后生成，勿手改（改源 MD 后重跑）。

\begin{gknotice}
\textbf{使用说明}：本节只给最终答案便于快速核对，完整推导见解析版。答案为编写者独立推导，\textbf{非官方标准答案}。
\end{gknotice}

\gkpart{选择题}
\begin{center}\begin{tabular}{ccccccccccccc}
\toprule
题号 & 1 & 2 & 3 & 4 & 5 & 6 & 7 & 8 & 9 & 10 & 11 & 12 \\
\midrule
答案 & C & A & C & C & A & D & B & B & D & B & C & C \\
\bottomrule
\end{tabular}\end{center}

\gkpart{填空题}
\begin{itemize}[leftmargin=2.6em]
  \item[\textbf{13.}] $P(\xi=0)=\dfrac1{10}$，$P(\xi=1)=\dfrac35$，$P(\xi=2)=\dfrac3{10}$
  \item[\textbf{14.}] $5$
  \item[\textbf{15.}] $\dfrac{x^2}{2}+y^2=1$
  \item[\textbf{16.}] ①②④
\end{itemize}

\gkpart{解答题}
\begin{itemize}[leftmargin=2.6em]
  \item[\textbf{17.}] (1) $\tan A=2\tan B$；(2) $AB$ 边上的高为 $2+\sqrt6$
  \item[\textbf{18.}] (1) $\dfrac67$；(2) $\dfrac12$
  \item[\textbf{19.}] (1) $\left\{\dfrac{S_n}{n}\right\}$ 是首项为 $1$、公比为 $2$ 的等比数列；(2) $S_{n+1}=4a_n$
  \item[\textbf{20.}] (1) $CD\perp$ 平面 $BDM$；(2) $90^\circ$
  \item[\textbf{21.}] (1) $\arccos\left(-\dfrac3{\sqrt{41}}\right)$；(2) $l$ 在 $y$ 轴上的截距变化范围为 $\left[-\dfrac43,-\dfrac34\right]\cup\left[\dfrac34,\dfrac43\right]$
  \item[\textbf{22.}] (1) $f(x)$ 的最大值为 $0$；(2) 对 $t=b/a>1$，有 $0<t\ln t-(1+t)\ln\dfrac{1+t}{2}<(t-1)\ln2$，故原不等式成立
\end{itemize}
